\section{ML Pipeline}
\label{sec:ml}

\subsection{Isolation Forest}
Isolation Forest partitions the feature space using random axis-aligned cuts.
Anomalies require fewer partitions to isolate and therefore have shorter
average path lengths.
We convert raw decision-function scores to [0, 1] by min-max normalisation:
\begin{equation}
  s_{\text{IF}}(x) = \frac{f(x) - \max f}{\min f - \max f}
  \label{eq:iforest_norm}
\end{equation}
where $f(\cdot)$ is the raw decision function.
We set the contamination parameter to 5\% based on the observed attack
ratio in the training splits of both datasets.

\subsection{LSTM Autoencoder}
The autoencoder uses a sequence-to-sequence architecture with two LSTM
layers in the encoder (hidden sizes 64 and 32) and a mirrored decoder.
Input windows of $T = 10$ consecutive flows are reconstructed; the
reconstruction error is:
\begin{equation}
  e_t = \frac{1}{T \cdot d} \sum_{i=0}^{T-1} \| x_{t+i} - \hat{x}_{t+i} \|^2
  \label{eq:recon_error}
\end{equation}
The anomaly threshold $\tau$ is set at the 95th percentile of $e_t$ on
the training set.
Scores are normalised to [0, 1] by dividing by $\tau$.

\subsection{Weighted Ensemble}
The final threat score is:
\begin{equation}
  s(x) = 0.4 \cdot s_{\text{IF}}(x) + 0.6 \cdot s_{\text{AE}}(x)
  \label{eq:ensemble}
\end{equation}
A flow is classified as anomalous if $s(x) \geq 0.5$.
The weights were selected by grid search over
$\{(w_1, w_2) : w_1 + w_2 = 1,\; w_1 \in \{0.1, 0.2, \ldots, 0.9\}\}$
on the UNSW-NB15 validation split, maximising F1.
The autoencoder receives higher weight because temporal context provides
a stronger signal for the slow attacks (beaconing, exfiltration) that
dominate our threat model.

\subsection{Training Protocol}
Models are trained on the combined UNSW-NB15 training split (normal flows
only; contamination is handled by the iForest hyperparameter).
Training runs for up to 50 epochs with early stopping (patience = 5) on
validation MSE.
Both models are persisted to a Kubernetes PersistentVolume and mounted
read-only into the scoring API pod.
